\documentclass[11pt,a4paper]{article}
% Shared preamble of the FEBAUTO Racing manuals.
\usepackage[margin=2.2cm,headheight=16pt]{geometry}
\usepackage[T1]{fontenc}
\usepackage{lmodern}
\usepackage{xcolor}
\usepackage{graphicx}
\usepackage{booktabs}
\usepackage{tabularx}
\usepackage{listings}
\usepackage{fancyhdr}
\usepackage{enumitem}
\usepackage{tikz}
\usetikzlibrary{arrows.meta}
\usepackage[hidelinks,colorlinks=true,linkcolor=febblue,urlcolor=febblue]{hyperref}

\definecolor{febblue}{HTML}{1C4687}
\definecolor{febgold}{HTML}{FDB515}
\definecolor{codebg}{HTML}{F3F4F6}

\lstset{
  basicstyle=\ttfamily\small,
  backgroundcolor=\color{codebg},
  frame=single, framerule=0pt, framesep=4pt,
  xleftmargin=4pt, xrightmargin=4pt,
  breaklines=true, columns=fullflexible, keepspaces=true,
  showstringspaces=false,
}
\newcommand{\file}[1]{\texttt{#1}}
\newcommand{\cmd}[1]{\texttt{#1}}
\newcommand{\rulebox}[1]{\par\noindent\colorbox{codebg}{\parbox{\dimexpr\linewidth-2\fboxsep}{#1}}\par\medskip}

\pagestyle{fancy}
\fancyhf{}
\fancyhead[L]{\color{febblue}\textbf{FEBAUTO Racing}}
\fancyhead[R]{\color{febblue}\nouppercase{\leftmark}}
\fancyfoot[C]{\thepage}
\renewcommand{\headrulewidth}{0.8pt}
\renewcommand{\headrule}{\color{febgold}\hrule width\headwidth height\headrulewidth}
\setcounter{tocdepth}{1}

% \febtitle{subtitle}{part lines}
\newcommand{\febtitle}[2]{%
\begin{titlepage}
  \centering
  \vspace*{2cm}
  \includegraphics[width=5cm]{../../assets/feb_logo_512.png}\\[1.5cm]
  {\color{febblue}\Huge\bfseries FEBAUTO Racing}\\[6pt]
  {\Large #1}\\[2cm]
  \begin{tabular}{rl}#2\end{tabular}\\[3cm]
  Formula Electric at Berkeley\\
  September 2026, platform release v0.1.0\\[6pt]
  \url{https://github.com/FEBAutonomous/feb-racing}\\
  \url{https://febautonomous.github.io/feb-racing/}
\end{titlepage}
\tableofcontents
\newpage}

\setlist{nosep,leftmargin=*}

\begin{document}
\febtitle{feb\_racer: the memory racer}{%
  \textbf{What} & a lap to map the track, then a raceline driven by pure pursuit\\[3pt]
  \textbf{Where from} & the RoboRacer League~6 entry's mapping and raceline, on the car's contract\\[3pt]
  \textbf{Status} & scored practice 2026-10-24: loop 138.8~s and porto 110.8~s, clean, against the gap follower's 173.4 and 168.2; car: untried\\}

\section{What it does}

Three phases, visible in \cmd{./feb logs}:

\begin{tabularx}{\linewidth}{lX}
\toprule
Phase & What happens \\
\midrule
WARMUP & the gap follower drives one lap at \texttt{warmup\_speed} while a scan matcher builds an occupancy map
  in the start frame, correcting the odometry as the map fills. The lap closes when the car is back within
  \texttt{close\_radius\_m} of where it started, heading the same way, after \texttt{min\_lap\_m}. \\
COMPUTE & the map is frozen. A thread extracts the corridor centre, optimises a minimum-curvature line inside the
  corridor and gives it a friction-ellipse speed profile (\texttt{v\_max}, \texttt{ay\_max}, \texttt{a\_acc},
  \texttt{a\_dec}). Under a second; the car keeps driving reactively meanwhile. \\
RACE & every scan is matched against the map. A pure-pursuit tracker aims at the raceline point one lookahead ahead
  and commands the profile's speed over the next \texttt{speed\_horizon\_m}. The gap follower's clearance brake stays
  armed underneath. \\
\bottomrule
\end{tabularx}

\medskip
\textbf{When the map match drops} below \texttt{lost\_score} for \texttt{lost\_handoff\_s}, the car hands back to the
gap follower at \texttt{recover\_speed} and takes the line back when the map matches again. The reactive law has no
pose to lose, which is why it is the safety net; that idea comes straight from the league entry. After two seconds
lost, or whenever the dead-man re-engages, it searches the whole map for itself once.

\section{What it is not}

It is not the league entry copied over. That node was 850 lines tied to the league's throttle commands and encoder
ticks, with a sampling MPC and a dozen detectors for the league simulator's quirks. Kept: the scan matcher
(\file{localizer.py}) and the raceline generator (\file{raceline\_gen.py}), unchanged. Replaced: the controller, by
pure pursuit on the contract's speed and steering angle; the inputs, by \cmd{/scan} and \cmd{/odom}; the detectors,
by the dead-man. The racer is \file{racer.py}, under 300 lines.

\section{Running it}

\begin{lstlisting}
./feb up --driver feb_racer                         # simulator, it drives from the start
./feb up --deadman --driver feb_racer               # hold SPACE: warm-up, then it races; let go, it stops
./feb practice --driver feb_racer                   # scored: the warm-up lap is the rules' warm-up lap
./feb up --car --driver feb_racer --mode-name inspection   # the car: RB held, slow numbers from driver.car.yaml
\end{lstlisting}

The map lives in the node for as long as it runs. Stop and start the stack and it maps again.

\section{The numbers}

\file{config/driver.yaml}, every value commented. The ones to touch first:

\begin{tabularx}{\linewidth}{lX}
\toprule
Value & Meaning \\
\midrule
\texttt{warmup\_speed} & the mapping lap; slow is fine, the map only needs to be complete \\
\texttt{v\_max}, \texttt{ay\_max} & the speed profile's cap and lateral limit; the first thing to raise on a track it laps cleanly \\
\texttt{speed\_scale} & scales the whole profile; 0.8 for a first run on a new track \\
\texttt{wall\_margin\_m} & how far the line keeps from the walls on top of half a car \\
\texttt{lookahead\_gain} & pure pursuit's lookahead per m/s; higher is smoother and cuts corners more \\
\texttt{lost\_score} & below this the pose is not trusted \\
\bottomrule
\end{tabularx}

\medskip
\file{config/driver.car.yaml} is loaded on top on the car: slower warm-up, \texttt{v\_max} 2.5, scale 0.8, and a
longer \texttt{range\_max} because the S2 sees further than the simulator's lidar.

\section{What the scored runs taught}

The free runs were clean from the first evening; the scored practice, at the league's 10~Hz, was not,
twice. Both causes were measured, not guessed. First, our odometry's heading integrates the gyro and at
10~Hz drifts 18 degrees in 30~s and 84 by the end of a run; the racer now takes its heading from the
IMU's absolute orientation when there is one. Second, with a good pose the car still clipped walls:
at 10~Hz and a quarter-second actuator delay, pure pursuit at 4~m/s steers late. The tracker now
aims from where the car will be one delay later, keeps 0.25~m from the walls and caps at 3~m/s. The
same heading drift is why the cone racer lost localisation at 10~Hz; the house adapter will get the
IMU heading too.

\section{Known limits}

\begin{itemize}
  \item Walled tracks only. Cones are a different problem; see \texttt{feb\_cone\_racer}.
  \item A respawn or a carried car breaks the map frame until the global search finds it again.
  \item The car has not run it. The matcher and the map were built in the simulator with its lidar; the S2 is denser
    and sees further, which only helps, but nobody has watched it.
  \item Pure pursuit is a looser tracker than the league's MPC; the lap times have room.
\end{itemize}

\end{document}
